\subsubsection{Diffusion on the surface}

Now that we are working with simulations, we can make our model slightly more complex by dropping one of our assumptions. Let $\Omega$ and $\Gamma$ be constructed as in \cref{sec:bio_motiv}. In this chapter, we restrict our study to the dimensions of $d = 2$ and $d = 3$, since these are the most relevant to the biological applications of this model. Recall that in \cref{sec:bio_motiv}, we assumed the receptors, both free and bounded, to be fixed in place on the cell membrane. This resulted in a system of one PDE and two ODEs, the latter being both defined separately for every point on the cell membrane. We now drop this assumption and add diffusion of free and bounded receptors, leading to a new system of reaction diffusion PDEs given by

\begin{equation}
    \label{eq:threespeciesmodel}
    \begin{aligned}
    u_t &= D_u \Delta u + P_u + \alpha u_b - \delta_u u - 2 \kon u^2 v + 2 \koff u_b & \text{on } &(0, T) \times \Gamma, \\
    u_{b,t} &= D_{u_b} \Delta u_b - \delta_{u_b} u_b + \kon u^2 v - \koff u_b & \text{on } &(0, T) \times \Gamma, \\
    v_t &= D_v \Delta v - \delta_v v & \text{in } &(0, T) \times \Omega, \\
    D_v \pderiv[v]{\nu} &= P_v - \kon u^2 v + \koff u_b & \text{on } &(0, T) \times \Gamma, \\
    D_v \pderiv[v]{\nu} &= 0 & \text{on } &(0, T) \times \Gamma_e, \\
\end{aligned}
\end{equation}

\noindent
where $D_u$ and $D_{u_b}$ are positive. The corresponding initial conditions are

\begin{equation*}
\begin{aligned}
    u(0, x) = u_0(x), \qquad u_b(0, x) = u_{b,0}(x) \qquad \text{and} \qquad v(0, y) = v_0(y)
\end{aligned}
\end{equation*}

\noindent
for all $x \in \Gamma$ and $y \in \overline{\Omega}$.

There are multiple things to note about \cref{eq:threespeciesmodel}. Since $\Gamma$ has an empty interior, it is impossible to define spatial derivatives for functions living exclusively on this set. This means we would not be able to use the Laplace operator to include diffusion of the receptors as we do in \cref{eq:threespeciesmodel}. This issue is solved by using the so called the Laplace-Beltrami operator instead, which is a generalisation of the Laplace operator to Riemannian manifolds. Since we assumed $\Gamma$ to be a $C^1$ boundary, it forms such a manifold. The Laplace-Beltrami operator is further discussed in \cite{jost2017riemannian}, but we will not go into details here. For our purposes, it is sufficient to note that it behaves similarly to the classical Laplace operator and admits analogous rules of calculus such as partial integration. It is sometimes denoted by $\Delta_\Gamma$, but for the sake of simplicity, we will omit this notation and use $\Delta$ for both the Laplace and Laplace-Beltrami operators, but they are in fact different operators.

Since we now have PDEs for $u$ and $u_b$, it seems necessary to impose certain boundary conditions. However, recalling our exact construction of the domain, we can see that the set $\Gamma$ representing the cell membranes, when viewed as a separate $(d - 1)$-dimensional manifold, does not have a boundary. Since the equations for $u$ and $u_b$ strictly live on this manifold, boundary conditions are not required.

\subsubsection{Derivation of the two-species model}

We seek to compare this three-species bulk-surface model to a two-species model similar to that derived in $\cite{iber2014}$. To obtain such a two-species model, we make some of the same assumptions. The first and most impactful is that the dynamics of the bounded receptors are sufficiently fast to make $u_{b, t}$ and $\Delta u_b$ are negligible. This results in an explicit expression for $u_b$ in terms of $u$ and $v$,

\begin{equation*}
    u_b = \frac{\kon}{\delta_{u_b} + \koff} u^2 v .
\end{equation*}

\noindent
This is called a \emph{quasi-steady state}. This expression enables us to write a two-species model

\begin{equation}
\label{eq:twospeciesmodel}
\begin{aligned}
    u_t &= D_u \Delta u + P_u - \delta_u u + \beta u^2 v & \text{on } &(0, T) \times \Gamma, \\
    v_t &= D_v \Delta v - \delta_v v & \text{in } &(0, T) \times \Omega, \\
    D_v \pderiv[v]{\nu} &= P_v - \beta u^2 v & \text{on } &(0, T) \times \Gamma, \\
    D_v \pderiv[v]{\nu} &= 0 & \text{on } &(0, T) \times \Gamma_e, \\
\end{aligned}
\end{equation}

\noindent
where $\beta := \frac{\kon \delta_{u_b}}{\delta_{u_b} + \koff}$ and we make the second assumption, namely that $\alpha = 3\beta$. The final assumption made in \cite{iber2014} is that $\delta_v v \ll \beta u^2 v$, making the natural decay of the ligands negligible and recovering the classical Schnakenberg equations. However, as we saw in \cref{sec:schnakenberg_decay_example}, the Schnakenberg equations with additional decay of $v$ is also able to generate patterns, so we will not be making this assumption.